% Part:sets-functions-relations
% Chapter: size-of-sets
% Section: zig-zag

\documentclass[../../../include/open-logic-section]{subfiles}

\begin{document}

\olfileid[ar]{sfr}{siz}{zigzag}

\olsection{طريقة كانتور المتعرجة}

\begin{explain}
سبق أن نظرنا في بعض التعدادات «السهلة». والآن سننظر في أمر أصعب
قليلًا. تأمل مجموعة أزواج الأعداد
الطبيعية\oliflabeldef{sfr:set:pai:sec}{، التي عرّفناها في
\olref[set][pai]{sec} على النحو الآتي:}{ المعرّفة كما يأتي:}
\[
\Nat \times \Nat = \Setabs{\tuple{n,m}}{n,m \in \Nat}
\]
يمكننا تنظيم هذه الأزواج المرتبة في \emph{مصفوفة} كما يأتي:
\[
\begin{array}{ c | c | c | c | c | c}
& \mathbf 0 & \mathbf 1 & \mathbf 2 & \mathbf 3 & \dots \\
\hline
\mathbf 0 & \tuple{0,0} & \tuple{0,1} & \tuple{0,2} & \tuple{0,3} & \dots \\
\hline
\mathbf 1 & \tuple{1,0} & \tuple{1,1} & \tuple{1,2} & \tuple{1,3} & \dots \\
\hline
\mathbf 2 & \tuple{2,0} & \tuple{2,1} & \tuple{2,2} & \tuple{2,3} & \dots \\
\hline
\mathbf 3 & \tuple{3,0} & \tuple{3,1} & \tuple{3,2} & \tuple{3,3} & \dots \\
\hline
\vdots & \vdots & \vdots & \vdots & \vdots & \ddots\\
\end{array}
\]
من الواضح أن كل زوج مرتب في $\Nat \times \Nat$ سيظهر مرة واحدة
بالضبط في المصفوفة. وبخاصة، سيظهر $\tuple{n,m}$ في الصف ذي الرتبة
$n$ والعمود ذي الرتبة $m$. ولكن كيف ننظم عناصر مثل هذه المصفوفة في
قائمة «أحادية البعد»؟ يبين النمط في المصفوفة أدناه طريقةً لفعل ذلك
(مع أن ثمة، بطبيعة الحال، خيارات أخرى كثيرة):
\[
\begin{array}{ c | c | c | c | c | c | c}
& \mathbf 0 & \mathbf 1 & \mathbf 2 & \mathbf 3 & \mathbf 4 &\dots \\
\hline
\mathbf 0 & 0  & 1& 3 & 6& 10 &\ldots \\
\hline
\mathbf 1 &2 & 4& 7 & 11 & \dots &\ldots \\
\hline
\mathbf 2 & 5 & 8 & 12 & \ldots & \dots&\ldots \\
\hline
\mathbf 3 & 9 & 13 & \ldots & \ldots & \dots & \ldots \\
\hline
\mathbf 4 & 14 & \ldots & \ldots & \ldots & \dots & \ldots \\
\hline
\vdots & \vdots & \vdots & \vdots & \vdots&\ldots & \ddots\\
\end{array}
\]\noindent
يسمى هذا النمط \emph{طريقة كانتور المتعرجة}. وهو يعدّد
$\Nat \times \Nat$ كما يأتي:
\[
\tuple{0,0}, \tuple{0,1}, \tuple{1,0}, \tuple{0,2}, \tuple{1,1},
\tuple{2,0}, \tuple{0,3}, \tuple{1,2}, \tuple{2,1}, \tuple{3,0}, \dots
\]
وهذا يثبت النتيجة الآتية:
\end{explain}

\begin{prop}\ollabel{natsquaredenumerable}
تكون $\Nat \times \Nat$ !!{enumerable}.
\end{prop}

\begin{proof}
لتكن $f \colon \Nat \to \Nat\times\Nat$ دالة ترسل كل $k \in \Nat$
إلى الزوج $\tuple{n,m} \in \Nat \times \Nat$ الذي يحقق أن $k$ هو
القيمة المدرجة في الصف ذي الرتبة $n$ والعمود ذي الرتبة $m$ في مصفوفة
كانتور المتعرجة.
\end{proof}

\begin{explain}
وتقبل هذه التقنية أيضًا تعميمًا أنيقًا. فمثلًا، يمكننا استعمالها لتعداد
مجموعة الثلاثيات المرتبة من الأعداد الطبيعية، أي:
\[
\Nat \times \Nat \times \Nat = \Setabs{\tuple{n,m,k}}{n,m,k \in \Nat}
\]
ننظر إلى $\Nat \times \Nat \times \Nat$ بوصفها الضرب الديكارتي
لـ$\Nat \times \Nat$ في $\Nat$، أي
\[
\Nat^3 = (\Nat \times \Nat) \times \Nat =
\Setabs{\tuple{\tuple{n,m},k}}{n, m, k
  \in \Nat }
\]
ومن ثم يمكننا تعداد $\Nat^3$ بمصفوفة، فنضع تعداد $\Nat$ على أحد
محوريها، وتعداد $\Nat^2$ على المحور الآخر:
\[
\begin{array}{ c | c | c | c | c | c}
& \mathbf 0 & \mathbf 1 & \mathbf 2 & \mathbf 3 & \dots \\
\hline
\mathbf{\tuple{0,0}} & \tuple{0,0,0} & \tuple{0,0,1} & \tuple{0,0,2} & \tuple{0,0,3} & \dots \\
\hline
\mathbf{\tuple{0,1}} & \tuple{0,1,0} & \tuple{0,1,1} & \tuple{0,1,2} & \tuple{0,1,3} & \dots \\
\hline
\mathbf{\tuple{1,0}} & \tuple{1,0,0} & \tuple{1,0,1} & \tuple{1,0,2} & \tuple{1,0,3} & \dots \\
\hline
\mathbf{\tuple{0,2}} & \tuple{0,2,0} & \tuple{0,2,1} & \tuple{0,2,2} & \tuple{0,2,3} & \dots\\
\hline
\vdots & \vdots & \vdots & \vdots & \vdots & \ddots \\
\end{array}
\]
وبذلك، باستعمال طريقة شبيهة بطريقة كانتور المتعرجة، يمكننا أن نحصل
أيضًا على تعداد لـ$\Nat^3$. ويمكننا مواصلة ذلك، فنحصل على تعدادات
لـ$\Nat^n$ لكل عدد طبيعي $n$. ومن ثم لدينا:
\end{explain}

\begin{prop}
تكون $\Nat^n$ !!{enumerable}، لكل $n \in \Nat$.
\end{prop}

\begin{prob}\label{sfr:siz:zigzag:prob:posint-n}
أثبت أن $(\PosInt)^n$ !!{enumerable}، لكل $n \in \Nat$.
\end{prob}

\begin{prob}\label{sfr:siz:zigzag:prob:posint-star}
  أثبت أن $(\PosInt)^*$ !!{enumerable}. ويجوز لك افتراض \cref{sfr:siz:zigzag:prob:posint-n}.
\end{prob}

\end{document}
